
\documentclass{article}
\usepackage{colortbl}
\usepackage{makecell}
\usepackage{multirow}
\usepackage{supertabular}

\begin{document}

\newcounter{utterance}

\centering \large Interaction Transcript for game `ifeval', experiment `ifeval', episode 136 with Bentel rockers\_\_qwen3.5{-}4b\_\_Bentel\_iter\_2.
\vspace{24pt}

{ \footnotesize  \setcounter{utterance}{1}
\setlength{\tabcolsep}{0pt}
\begin{supertabular}{c@{$\;$}|p{.15\linewidth}@{}p{.15\linewidth}p{.15\linewidth}p{.15\linewidth}p{.15\linewidth}p{.15\linewidth}}
    \# & \multicolumn{2}{c}{Player} && \multicolumn{2}{c}{Game Master} \\
    \hline

    \theutterance \stepcounter{utterance}  
    & & & \multicolumn{4}{p{0.6\linewidth}}{
        \cellcolor[rgb]{0.9,0.9,0.9}{
            \makecell[{{p{\linewidth}}}]{
                \texttt{\tiny{[P1$\langle$GM]}}
                \texttt{Write an interesting and funny article about the biology of a banana peel. In your response, the word disappointed should appear at least 2 times, and at least six section should be highlighted with markdown,  i.e *banana peel*.} \\
            }
        }
    }
     \\ \\

    \theutterance \stepcounter{utterance}  
    & \multicolumn{4}{p{0.6\linewidth}}{
        \cellcolor[rgb]{0.561,0.737,0.561}{
            \makecell[{{p{\linewidth}}}]{
                \texttt{\tiny{[P1$\rangle$GM]}}
                \texttt{\# The Biology of a Banana Peel} \\
\\ 
\texttt{\#\# Introduction} \\
\texttt{Bananas are a popular fruit that is often enjoyed for their sweet taste and creamy texture. However, the banana peel is an interesting part of the fruit that is often overlooked. In this article, we will explore the biology of the banana peel and its role in the overall structure of the fruit.} \\
\\ 
\texttt{\#\# The Structure of the Banana Peel} \\
\texttt{The banana peel is a protective layer that surrounds the fruit. It is made up of a thin layer of skin and a thicker layer of flesh. The skin is typically green when the fruit is unripe and turns yellow when it is ripe. The flesh is a soft, creamy tissue that is rich in nutrients.} \\
\\ 
\texttt{\#\# The Function of the Banana Peel} \\
\texttt{The banana peel serves several important functions in the life of the fruit. It protects the fruit from physical damage and helps to regulate the exchange of gases between the fruit and the surrounding environment. The peel also contains enzymes that help to break down the fruit and release its nutrients.} \\
\\ 
\texttt{\#\# The Development of the Banana Peel} \\
\texttt{The banana peel begins to form when the fruit is still a small, green} \\
            }
        }
    }
    & & \\ \\

    \theutterance \stepcounter{utterance}  
    & & & \multicolumn{2}{p{0.3\linewidth}}{
        \cellcolor[rgb]{0.9,0.9,0.9}{
            \makecell[{{p{\linewidth}}}]{
                \texttt{\tiny{[GM$|$GM]}}
                {'keywords:frequency': {'relation': 'at least', 'keyword': 'disappointed', 'frequency': 2}, 'detectable_format:number_highlighted_sections': {'num_highlights': 6}}
            }
        }
    }
    & & \\ \\

    \theutterance \stepcounter{utterance}  
    & & & \multicolumn{2}{p{0.3\linewidth}}{
        \cellcolor[rgb]{0.9,0.9,0.9}{
            \makecell[{{p{\linewidth}}}]{
                \texttt{\tiny{[GM$|$GM]}}
                \texttt{game\_result = LOSE} \\
            }
        }
    }
    & & \\ \\

\end{supertabular}
}

\end{document}
